% Non-Latin script + symbol font handling for xelatex.
% TEMPLATE FILE — processed by build_book.py at build time.
% The token Adobe Devanagari is replaced with the value of `devanagarifont`
% from as_book.yaml; the rendered file is written to build/devanagari-preamble.tex
% and included via pandoc's --include-in-header flag.
%
% Strategy: explicit per-run font wrapping (no ucharclasses).
%   - mainfont (set via as_book.yaml) handles Latin + IAST diacritics.
%   - build_book.py's assemble step scans for runs of each script's Unicode
%     range and wraps them in raw-LaTeX `{\fontname …}` so the font switch
%     is unconditional inside the wrap group and reverts cleanly at the
%     closing brace.
%   - Each script declared here covers one Unicode block (or related blocks);
%     the build-script wrapper lists the exact range per script.
%   - ucharclasses was tried earlier and removed (it routes some punctuation
%     codepoints through a separate class that overrides the wrap-group's
%     font selection).

\usepackage{fontspec}
\usepackage{caption}
\usepackage[normalem]{ulem}

% Figure numbers are authored directly in Markdown captions (for example,
% Figure 0.1 and Figure E.3). Suppress LaTeX's independent sequential label so
% numbered and intentionally unnumbered captions both render as written.
\captionsetup[figure]{labelformat=empty}

% Keep ordinary figures beside the surrounding prose whenever space permits.
% Reserve float-only pages for figures that occupy most of the text area.
\renewcommand{\topfraction}{0.9}
\renewcommand{\bottomfraction}{0.8}
\renewcommand{\textfraction}{0.1}
\renewcommand{\floatpagefraction}{0.85}

\newfontfamily{\devanagarifont}[Script=Devanagari]{Adobe Devanagari}
% Adobe Devanagari does not cover the Vedic Extensions block. Use a Sanskrit
% face only for those signs so ordinary Devanagari typography remains intact.
\newfontfamily{\vedicfont}[Script=Devanagari]{Tiro Devanagari Sanskrit}
\newfontfamily{\arabicfont}[Script=Arabic]{Geeza Pro}
\newfontfamily{\hebrewfont}[Script=Hebrew]{Arial Hebrew}
\newfontfamily{\tamilfont}[Script=Tamil]{Tamil Sangam MN}
\newfontfamily{\telugufont}[Script=Telugu]{Telugu Sangam MN}
\newfontfamily{\kannadafont}[Script=Kannada]{Kannada Sangam MN}
\newfontfamily{\malayalamfont}[Script=Malayalam]{Malayalam Sangam MN}
\newfontfamily{\jpfont}{Hiragino Sans}
\newfontfamily{\cjkfont}{PingFang SC}
\newfontfamily{\oldpersianfont}{Noto Sans Old Persian}
\newfontfamily{\gothicfont}{Noto Sans Gothic}
\newfontfamily{\avestanfont}{Noto Sans Avestan}
% Symbol / Latin-Extended fallback for arrows, subscripts, modifier letters,
% IAST diacritics Charter lacks (ē, ʷ, ʾ, ḱ, ẓ, →, ←, ₁, ₂, μ, τ, ρ, etc.).
% Arial Unicode MS chosen for breadth of Unicode coverage and reliable
% presence on macOS systems. (STIX Two Text was tried but lacks U+2192 arrow.
% Noto Sans was tried but not all macOS systems ship the un-suffixed family
% — only script-specific variants like "Noto Sans Devanagari" are present.)
\newfontfamily{\symbolfont}{Arial Unicode MS}

% --- Sanskrit chapter-opening epigraph typography ---
%
% Deploy inside a markdown blockquote, on its own line after the Devanagari
% verse and the IAST gloss. The macro right-aligns the attribution with an
% em-dash prefix and italicizes the source name. Pair with a \bigskip raw-LaTeX
% block after the blockquote for vertical separation from the chapter body.
%
% Markdown form:
%   > ऋटुरषाणां मूर्धा ।
%   >
%   > *ṛṭuraṣāṇāṃ mūrdhā.*
%   >
%   > \episource{Pāṇinīya Śikṣā tradition}
%
%   \bigskip
%
% Resulting LaTeX: \hfill{\itshape\textemdash\ Pāṇinīya Śikṣā tradition}
\newcommand{\episource}[1]{\hfill{\itshape\textemdash\ #1}}

% Epigraph environment. Pairs with the markdown fenced-div wrapper
% `::: epigraph ... :::` that wraps the verse + IAST + \episource{} lines.
% Applies a slightly larger font size for prominence; leaves \episource{}
% (right-aligned italic em-dash + source) and the IAST *...* markdown
% italics to handle their own typographic work.
\newenvironment{epigraph}{\par\medskip\bgroup\large}{\egroup\par\medskip}

% --- \partopener: part title + opener prose on the same page ---
%
% Modified \part that allows the part opener prose (the framing paragraph
% under the part title) to flow on the part-title page rather than being
% pushed to the next page by \cleardoublepage.
%
% Standard book-class \part centers the title vertically on a near-empty
% page and forces \cleardoublepage afterward, so opener prose lands on
% the page following the title page. \partopener instead:
%   - Forces a new page (the part still starts on a new recto)
%   - Places the title near the top of the page (top vskip), not centered
%   - Suppresses the trailing \cleardoublepage so subsequent document
%     content (the opener prose) flows on the same page
%
% Usage:
%   \partopener[TOC TEXT]{PAGE TITLE}
%   <opener prose flows as ordinary document content here>
%
% If TOC TEXT is omitted, PAGE TITLE is used for the TOC entry.
\makeatletter
\newcommand{\partopener}[2][]{%
  \if@openright\cleardoublepage\else\clearpage\fi
  \thispagestyle{plain}%
  \refstepcounter{part}%
  \def\partopener@toc{#1}%
  \ifx\partopener@toc\@empty
    \addcontentsline{toc}{part}{#2}%
  \else
    \addcontentsline{toc}{part}{#1}%
  \fi
  \markboth{}{}%
  \vspace*{4em}% top margin before the title
  {\centering
   \interlinepenalty\@M
   \normalfont
   \Huge\bfseries #2\par}%
  \vskip 3em\relax
  % document flow continues here — opener prose appears on this page
}
\makeatother

% Pre-warm each font at document start so the first wrap-group transition
% does not leave its first glyph in the outgoing font.
\AtBeginDocument{%
  \setbox0=\hbox{\devanagarifont क}%
  \setbox0=\hbox{\arabicfont ا}%
  \setbox0=\hbox{\hebrewfont א}%
  \setbox0=\hbox{\tamilfont அ}%
  \setbox0=\hbox{\telugufont అ}%
  \setbox0=\hbox{\kannadafont ಅ}%
  \setbox0=\hbox{\malayalamfont അ}%
  \setbox0=\hbox{\jpfont ア}%
  \setbox0=\hbox{\cjkfont 中}%
  \setbox0=\hbox{\gothicfont 𐌾}%
  \setbox0=\hbox{\symbolfont →}%
}
